\section{Introdução}
\label{introducao}

A comunicação acadêmica nas universidades brasileiras acontece de forma fragmentada. Na Universidade Federal de Itajubá (UNIFEI), o SIGAA e o Moodle convivem e se sobrepõem: os dois guardam material de aula, os dois organizam a entrega de tarefas e ambos oferecem alguma forma de fórum. Ainda assim, quando surge uma dúvida sobre o conteúdo de uma disciplina, a conversa migra para o WhatsApp ou para o e-mail. O problema, portanto, não é ausência de recurso. É que o recurso existe e não é usado, e a discussão termina em um canal onde se mistura com outros assuntos e some do alcance de quem venha a ter a mesma dúvida no semestre seguinte. Sobre esse movimento, \cite{kent2013} relata que os estudantes preferem as redes sociais informais em vez dos fóruns oficiais dos sistemas de gestão de aprendizagem.

Vale examinar por que isso acontece. Os fóruns desses sistemas foram construídos como um item a mais dentro de uma plataforma que existe para outra coisa — matrícula, nota, documento —, e herdaram a lógica administrativa do resto. As discussões não se organizam por disciplina de um jeito que o aluno reconheça, o que faz com que o conhecimento produzido em um semestre não fique disponível para os próximos. Outras universidades apresentam comportamento parecido, como a de São Paulo (USP), a Universidade Federal de Minas Gerais (UFMG) e a Universidade Estadual do Paraná (UNESPAR), que usam o Moodle ou sistemas próprios com fórum previsto, mas pouco frequentado. Seguindo essa linha de raciocínio, a revisão sistemática feita por \cite{almatrafi2019} sobre fóruns em cursos online mostra que esses espaços, quando bem integrados às disciplinas, auxiliam no engajamento dos estudantes, e \cite{kumi2023} acrescenta a essa discussão que, ao usar a análise de redes, verificou as interações dentro de um fórum, revelando o papel que cada aluno ocupa na troca de respostas.

Há ainda um terceiro espaço que fica de fora de todos esses sistemas. O trabalho em grupo ocupa boa parte da carga de várias disciplinas e é combinado inteiramente por fora: a turma se divide por mensagem, discute a repartição de tarefas em conversa privada e entrega o resultado sem que nada desse percurso passe pela universidade. Quando o grupo trava, a dúvida raramente chega ao professor. Levá-la até ele significaria expor a conversa inteira, e nenhum grupo faz isso. \cite{scager2016} observa que a aprendizagem colaborativa depende de condições criadas de propósito, e não apenas da reunião de alunos em torno de uma tarefa. O que se perde aqui não é só o registro da discussão, mas a chance de quem ensina perceber onde a turma está emperrada enquanto ainda há tempo de intervir.

O voluntariado universitário sofre de um mal parecido, com contornos próprios. Projetos de extensão, ONGs e a atlética divulgam suas atividades sobretudo nas redes sociais, o que restringe o alcance das oportunidades a quem acompanha na hora da publicação. Muitos alunos descobrem a vaga tarde demais, repassada por terceiros, ou depois que a inscrição já fechou. A UNIFEI tem o módulo de Ações de Extensão dentro do SIGAA, mas ele foi pensado para o servidor da instituição, e não para o aluno: a tela de busca traz mais de quinze filtros institucionais, como Área do CNPq, Centro da Ação, Edital, Programa Estratégico e Dimensão Acadêmica, parâmetros que fazem sentido para uma pró-reitoria e pouco ajudam quem apenas procura uma vaga de monitoria ou um projeto que combine com o seu curso. A participação segue igualmente manual. O coordenador envia uma lista à pró-reitoria de extensão, os certificados de participação são emitidos um a um, e cabe ao próprio aluno solicitar o aproveitamento dessas horas como atividade complementar no seu histórico. Sobre esse ponto, \cite{lopesjr2023} relata que não basta a universidade ofertar os projetos, sendo necessário um esforço intencional da instituição para que os alunos de fato participem das atividades de voluntariado.

Diante desses problemas, este trabalho desenvolveu uma plataforma web que reúne, em um único ambiente, um fórum acadêmico organizado por disciplina, um módulo de trabalhos em grupo e um sistema de voluntariado com certificação automática. No fórum, cada disciplina ganha o seu próprio espaço, e a pessoa enxerga apenas aquelas em que tem vínculo — o que evita que a dúvida de um oitavo período apareça para quem está no primeiro. Alunos, monitores e professores entram com perfis distintos e podem criar tópicos, responder, votar, marcar a melhor resposta e sinalizar quando uma resposta não resolveu, além de moderar o conteúdo das disciplinas em que atuam. A busca reconhece dúvidas equivalentes escritas com palavras diferentes. Esse detalhe responde diretamente ao problema do conhecimento que se perde: uma pergunta sobre laço que executa uma vez a mais e outra sobre erro de \textit{off-by-one} descrevem exatamente a mesma coisa sem compartilhar um único termo.

O módulo de trabalhos em grupo nasceu da lacuna descrita antes. O professor propõe a atividade, anexa o material de apoio e decide como as equipes se formam, por escolha dos alunos ou por sorteio. Cada grupo recebe uma conversa privada em tempo real, na qual o professor não entra. Quando a equipe emperra, ela marca a mensagem exata da dúvida e a encaminha à monitoria, que recebe aquele trecho e a descrição do impasse — nada mais. A solução preserva os dois lados da tensão: o grupo mantém a privacidade que o faz permanecer na plataforma em vez de migrar para aplicativos externos, e quem ensina passa a enxergar onde a turma está travando. No módulo de voluntariado, as organizações cadastram suas oportunidades, os alunos se inscrevem online e, ao encerrar a atividade, recebem automaticamente o certificado de participação, documento que instrui o pedido de aproveitamento das horas como atividade complementar e que hoje depende da rotina manual descrita anteriormente. Cada certificado carrega um código de validação pública, o que permite a qualquer instância verificar a sua autenticidade sem depender da organização emissora. Sustentando tudo isso, um sistema de notificações em tempo real avisa quando surge uma resposta em um tópico, uma mensagem no grupo ou a confirmação de uma inscrição.

O objetivo geral foi projetar, desenvolver e avaliar essa plataforma como resposta aos problemas relatados, com foco na comunidade acadêmica da UNIFEI. O percurso passou pelo levantamento bibliográfico nos eixos de comunicação em ambientes educacionais, fóruns de discussão, aprendizagem colaborativa e voluntariado universitário; pela análise comparativa entre as quatro instituições citadas; pela definição dos requisitos; pela especificação da arquitetura; pela modelagem do banco de dados em terceira forma normal; e pela implementação, que hoje conta com vinte e seis tabelas, cobertura automatizada de testes e as telas de todos os perfis em funcionamento. A documentação completa dos requisitos funcionais e não funcionais, junto ao código-fonte, está disponível no repositório do projeto\footnote{Disponível em: \url{https://github.com/kellyr5/plataforma-unifei}}, escolha que mantém este texto concentrado nas decisões de projeto em vez de reproduzir listagens extensas.

A metodologia é a de uma pesquisa aplicada de natureza qualitativa, conduzida pelo desenvolvimento experimental do sistema em abordagem incremental, com funcionalidades entregues em ciclos curtos e validadas antes do início do conjunto seguinte. As tecnologias foram organizadas por camadas: React.js, TypeScript e TailwindCSS na apresentação; Python com Django, Django REST Framework e Django Channels na aplicação; PostgreSQL como banco relacional e Redis para o gerenciamento de tokens e a comunicação em tempo real. O escopo foi delimitado como um produto mínimo apaixonante, denominado MLP (\textit{Minimum Lovable Product}), que reúne as funcionalidades essenciais e o cuidado com a experiência do usuário desde a primeira versão. A escolha se justifica por uma razão prática: a plataforma disputa espaço com canais informais já consolidados entre os alunos, como o WhatsApp, o e-mail e o Discord, e uma ferramenta institucional que exija esforço de adaptação sem oferecer ganho evidente tende a ser abandonada nas primeiras semanas \citep{kniberg2016}.

O restante do artigo segue assim: a Seção 2 apresenta os conceitos e trabalhos que fundamentam a proposta, com a análise comparativa entre as universidades; a Seção 3 detalha a arquitetura, os requisitos, a modelagem do banco de dados e os diagramas que representam o funcionamento da plataforma; a Seção 4 reúne os resultados obtidos com a implementação; e a Seção 5 traz as conclusões e os trabalhos futuros.
